\documentclass[8pt]{article}
\usepackage{graphicx} % Required for inserting images
\usepackage{amsthm}
\usepackage{amsmath}
\usepackage{amsfonts}
\usepackage{amsopn}
\usepackage{comment}
\usepackage{amssymb}
\usepackage{hyperref}
\usepackage{bussproofs}
\usepackage[all, 2cell]{xy}
\usepackage[all]{xy}
\usepackage{rotating}
\usepackage{lscape}
\usepackage{minted}
\usepackage{dsfont}
\usepackage{tikz}
\usepackage{tikz-cd}
\usepackage{stmaryrd}
\usepackage{multicol}
\usepackage{cmll}
\usepackage{wasysym}
\usetikzlibrary{shapes.geometric, arrows, positioning, backgrounds}
\usetikzlibrary{cd}
\makeatletter
\tikzcdset{
  eq node/.style={
    commutative diagrams/math mode=false, anchor=center},
  eq/.style={
    phantom,
    /tikz/every to/.append style={
      edge node={node[commutative diagrams/eq node]
        {\@eqnswtrue\make@display@tag\ltx@label{#1}}}}}}
\makeatother

\theoremstyle{definition}
\newtheorem{definition}{Definition}[section]

\theoremstyle{definition}
\newtheorem{theorem}{Theorem}[section]

\theoremstyle{definition}
\newtheorem{claim}{Claim}[section]

\theoremstyle{definition}
\newtheorem{ex}{Example}[section] 

\theoremstyle{definition}
\newtheorem{cons}{Construction}[section] 

\theoremstyle{definition}
\newtheorem{rem}{Remark}[section] 


\theoremstyle{definition}
\newtheorem{prop}{Proposition}[section]

\theoremstyle{definition}
\newtheorem{lemma}{Lemma}[section]

\theoremstyle{definition}
\newtheorem{fact}{Fact}[section]

\theoremstyle{definition}
\newtheorem{remark}{Remark}[section]

\theoremstyle{definition}
\newtheorem{notation}{Notation}[section]

\theoremstyle{definition}
\newtheorem{example}{Example}[section]

\theoremstyle{definition}
\newtheorem{exercise}{Exercise}[section]

\theoremstyle{definition}
\newtheorem{col}{Corollary}[section]

\theoremstyle{question}
\newtheorem{question}{Question}

\let\strokeL\L
\renewcommand\L{\mathbf{L}}

\DeclareMathOperator{\bang}{!}

\DeclareMathOperator*{\colim}{\operatorname{colim}}
\newcommand{\Ob}[1]{\operatorname{Ob}({\mathcal{#1}})}
\newcommand{\Cov}[1]{\operatorname{Cov}(#1)}

\newcommand{\circvee}{\mathbin{\text{\ooalign{\hfil$\vee$\hfil\cr\kern0.1ex$\bigcirc$}}}}
\newcommand{\circwedge}{\mathbin{\text{\ooalign{\hfil$\wedge$\hfil\cr\kern0.1ex$\bigcirc$}}}}

\title{Notes A Functorial Excursion Between Algebraic Geometry and Linear Logic}
\author{Daniel Rogozin}
\date{ }

\newcommand{\IntHom}[1]{[#1]}

\begin{document}

\maketitle

\tableofcontents

\section{Basics}

Let us start by recalling basic principles of algebraic geometry. The intuition is that a space $X$ comes along with a topological space
$(X,\Sigma_X)$ such that every open set $U$ is equipped with a set $\mathcal{O}_X(U)$ of local operations called \emph{regular functions}. 
Such regular functions are typically defined as polynomials with division and each  $\mathcal{O}_X(U)$ defines a commutative ring.
In functorial terms, every space gives rise to a ringed space, that is, a topological space $(X, \Sigma_X)$ equipped with a contravariant functor
$\mathcal{O}_X : \Sigma^{op}_X \to {\bf Ring}$ from the set of opens $\Sigma^{op}_X$ with the reversed inclusion order. The functorial action
of $\mathcal{O}_X$ describes the restriction functions 
\[
\mathcal{O}_X(V \subseteq U) : \mathcal{O}_X(U) \to \mathcal{O}_X(V)
\]
which takes an operation $f \in \mathcal{O}_X(U)$ defined on an open set $U$ and restrict it to an operation
\[
f|_V := \mathcal{O}_X(V \subseteq U)(f)
\]
on an open subset $V \subseteq U$. The functorial understanding of the concept of a space is supported by the observation
that every commutative ring $R$ defines a ringed space:
\[
\operatorname{Spec}R = (\operatorname{Spec}R, \mathcal{O}_{\operatorname{Spec}R})
\]
on the set of prime ideals $\mathfrak{p}$ of the commutative ring $R$. Such a ringed space is called the \emph{affine scheme} associated to $R$.
Every space $X$ of the theory, called a \emph{scheme}, is a patchwork of affine schemes $\operatorname{Spec}R_i$ glued together using the language of sheaf theory. 

Every scheme $X$ comes equipped with a symmetric monoidal category ${\bf qcMod}_X$ of quasicoherent sheaf of $\mathcal{O}_X$-modules.
By a presheaf of $\mathcal{O}_X$-modules $M$ of a ringed space $(X, \mathcal{O}_X)$, we mean a contravariant functor
\[
M : \Omega^{op}_X \to {\bf Ab}
\]
where ${\bf Ab}$ is the category of Abelian groups such that each Abelian group $M(U)$ defines a $\mathcal{O}_X(U)$-module.
The Serre-Swann theorem states a vector bundle on a suitable ringed space $(X, \mathcal{O}_X)$ can be equivalently encoded as its sheaf of sections
that defines a locally free and projective sheaf of $\mathcal{O}_X$-modules. A sheaf of $\mathcal{O}_X$ is called quasicoherent
when it is locally presentable, that is, it is locally the cokernel of a morphism of free modules. The category ${\bf qcMod}_X$ extends 
the category of vector bundles to define an Abelian category where kernels and direct images can be computed.

The idea is to associate a model of intuitionistic linear logic, where formulas are understood as generalised vector bundles and proofs are vector fields
by building an appropriate linear-non-linear adjunction. To associate a model of linear logic with each scheme $X$, we take the category
${\bf PshMod}_X$ of presheafs of $\mathcal{O}_X$-modules and their homomorphisms as the underlying symmetric monoidal closed category with tensor product $\otimes_X$,
the structural sheaf $\mathcal{O}_X$ as unit, and $\multimap_X$ for internal homs. For the cartesian part 
of a linear-non-linear adjunction, we take the category ${\bf PshCoAlg}_X$ of cocommutative comonoids (or commutative coalgebras) in ${\bf PshMod}_X$.
It is convenient to the approach based on the functor of points (see \cite[I.4]{eisenbud2000geometry}). $\operatorname{Spec}$
defines a contravariant functor $\operatorname{Spec} : {\bf Ring}^{op} \to {\bf Scheme}$ that induces the \emph{nerve functor}
\[
{\bf nerve} : {\bf Scheme} \to [{\bf Ring},{\bf Set}]
\]
associating to every commutative ring $R$ \emph{the set of $R$-points in a scheme $X$} defined as ${\bf Scheme}(\operatorname{Spec}R, X)$.
The nerve functor is fully faithful and it allows one to see every scheme $X$ as a covariant presheaf
\[
{\bf nerve}(X) : {\bf Ring} \to {\bf Set}
\]
So we move to schemes to ${\bf Ring}$-spaces and associate a model of intuitionistic linear logic with each ${\bf Ring}$-space $X$, 
where formulas are interpreted in ${\bf PshMod}_X$ as presheaves of $\mathcal{O}_X$-module understood as generalised vector bundles
on the ${\bf Ring}$-space. The linear-non-linear adjunction (and therefore the exponential modality) is constructed by the Sweedler construction.

\section{On the category of rings}

Let $(\mathcal{A}, \otimes, \mathds{1})$ be a symmetric monoidal category with reflexive coequalisers preserved by tensor products in every coordinate.
For example, the category ${\bf Ab}$ of Abelian groups and homomorphisms. A ring is a monoid object $(R, m, e)$ in $(\mathcal{A}, \otimes, \mathds{1})$.
In other words, a ring is a triple such that $(R, m, e)$ consisting of an object $R$ and morphisms $m : R \otimes R \to R$ and $e : \mathds{1} \to R$
such that the following diagrams commute:
\[
\begin{tikzcd}
  R \otimes R \otimes R \ar[d, "R \otimes m"'] \ar[rr, "m \otimes R"] && R \otimes R \ar[d, "m"]  & R \otimes R \ar[d, "m"'] \ar[rr, "m"] && R \ar[d, "e \otimes R"] \\
  R \otimes R \ar[rr, "m"'] && R                                                                  & R \ar[urr, equal] \ar[rr, "R \otimes e"'] && R \otimes R
\end{tikzcd}
\]

A ring $(R, m, e)$ is commutative is the following triangle commutes:
\[
\begin{tikzcd}
R \otimes R \ar[dr, "m"'] \ar[rr, "\gamma_{R,R}"] && R \otimes R \ar[dl, "m"] \\
& R
\end{tikzcd}
\]

A \emph{ring morphism} $u : (R, m_R, e_R) \to (S, m_S, e_S)$ is a morphism $u : R \to S$ in the underlying category $\mathcal{A}$ making the following diagrams commute:
\[
\begin{tikzcd}
  R \otimes R \ar[d, "m_S"'] \ar[rr, "u \otimes u"] && S \otimes S \ar[d, "m_S"] && \mathds{1} \ar[dr, "e_S"] \ar[dl, "e_R"'] \\
  R \ar[rr, "u"'] && S                                                 &   R \ar[rr, "u"'] && S
\end{tikzcd}
\]
The category ${\bf Ring}(\mathcal{A})$ of rings in $\mathcal{A}$ defined by the tensor product $R,S \mapsto R \otimes S$ in $\mathcal{A}$.

\section{On the category of $R$-modules and modules}

Given a ring $R$ in a symmetric monoidal category $\mathcal{A}$, an \emph{$R$-module} is a pair $(M, \operatorname{act})$ where $M$ is an object
in $\mathcal{A}$ along with a morphism $\operatorname{act} : R \otimes M \to M$ in $\mathcal{A}$ such that:
\[
\begin{tikzcd}
  R \otimes R \otimes M \ar[d, "R \otimes \operatorname{act}"'] \ar[rr, "m_R \otimes M"] && R \otimes M \ar[d, "\operatorname{act}"] && R \otimes M \ar[dr, "\operatorname{act}"] \\
  R \otimes M \ar[rr, "\operatorname{act}"'] && M & M \ar[ur, "e_R \otimes M"] \ar[rr, equal] && M
\end{tikzcd}
\]
Equivalently, an R-module is an Eilenberg-Moore algebra for the monad $A \mapsto R \otimes A$. An $R$-module homomorphism $f : M \to N$
between two modules is a map $f$ in $\mathcal{A}$ such that the following diagram commutes:
\[
\begin{tikzcd}
  R \otimes M \ar[rr, "R \otimes f"] \ar[d, "\operatorname{act}_M"'] && R \otimes N \ar[d, "\operatorname{act}_N"]\\
  M \ar[rr, "f"'] && N
\end{tikzcd}
\]

${\bf Mod}_R$ is the category of $R$-modules and their morphisms.

A \emph{module} is a pair $(R, M)$ where $R$ is a ring and $M$ is a module. A \emph{module homomorphism} $(u, f) : (R, M) \to (S, N)$
is a ring morphism $u : R \to S$ and a morphism $f : M \to N$ such that
\[
\begin{tikzcd}
  R \otimes M \ar[rr, "u \otimes f"] \ar[d, "\operatorname{act}_M"'] && S \otimes N \ar[d, "\operatorname{act}_N"] \\
  M \ar[rr, "f"'] && N
\end{tikzcd}
\]
${\bf Mod}$ is the category of modules and their morphisms.

There is a functor $\pi : {\bf Mod} \to {\bf Ring}$ transporting every module $(R, M)$ to the underlying ring $R$ and every morphism
$(u, f)$ to $u$. We will write
\[
u : R \to S \models f : M \to N
\]
for a module morphism $(u, f) : (R, M) \to (S, N)$. The intuition behind the notation is that every ring homomorphism
$u : R \to S$ induces a ``fibre'' consisting of all module morphisms $(u, f) : (R, M) \to (S, N)$ living above and refining the ring morphism
$u : R \to S$. The fibre of $\pi R$ for a commutative ring is the category ${\bf Mod}_R$. The functor $\pi : {\bf Mod} \to {\bf Ring}$ defines 
a Grothendieck fibration: every pair consisting of a ring morphism $u : R \to S$ and of an $S$-module $(N, \operatorname{act}_N)$ induces
an $R$-module noted 
\[{\bf res}_u N = (N, \operatorname{act}'_N)
\]
with the action morphism $\operatorname{act}'_N$ defined as the composite:
\[
R \otimes N \xrightarrow{u \otimes N} S \otimes N \xrightarrow{\operatorname{act}_N} N
\]
Moreover, the $S$-module with a module morphism:
\[
u : R \to S \models id_{N} : {\bf res}_u N \to N
\]
which is weakly Cartesian in the sense that every module morphism
\[
u : R \to S \models f : M \to N
\]
factors uniquely as 
\[
M \xrightarrow{(id_{R}, h)} {\bf res}_u N \xrightarrow{(u, id_N)} N
\]
for a morphism $h : M \to {\bf res}_u N$ in the category ${\bf Mod}_R$. The unique solution $h$ is equal to the original morphism $f : M \to N$
viewed as a morphism in ${\bf Mod}_R$. Given a ring morphism $u : R \to S$, the existence of a weakly Cartesian map for every $S$-module $N$
ensures the existence of a functor
\[
{\bf res}_u : {\bf Mod}_S \to {\bf Mod}_R
\]
called restriction of scalar along $u$. Along with the fact that the composite of weakly Cartesian maps is weakly Cartesian we have that
the forgetful functor $\pi : {\bf Mod} \to {\bf Ring}$ is a Grothendieck fibration.

Now we use that fact the category $\mathcal{A}$ has reflexive coequalisers, so one can show that ${\bf res}_u$ has a left adjoint
\[
{\bf ext}_u : {\bf Mod}_R \to {\bf Mod}_S
\]
constructed as follows. Every ring homomorphism $u : R \to S$ induces an $R \otimes S$-module notes as $R \otimes_u S$ defined as the reflexive 
coequaliser of the following diagram
\[
\begin{tikzcd}
  R \otimes R \otimes S \arrow[rrr, shift left=4, "m_R \otimes S"] \arrow[rrr, shift right=4, "(R \otimes m_S) \circ (R \otimes u \otimes S)"'] &&& R \otimes S \ar[lll, "R \otimes e_R \otimes S" description]
\end{tikzcd}
\]
computed in $\mathcal{A}$. Given three commutative rings $R$, $S_1$ and $S_2$, we define the functor 
\[
\circledast_R : {\bf Mod}_{S_1 \otimes R} \times {\bf Mod}_{R \otimes S_2} \to {\bf Mod}_{S_1 \otimes S_2}
\]
The two maps $\operatorname{act}_M : M \otimes R \to M$ and $\operatorname{act}_N : N \otimes R \to N$ are deduced by restriction of scalar
from the $S_1 \otimes R$-module structure of $M$ and the $R \otimes S_2$-module structure of $N$. The left adjunction functor ${\bf ext}_u$
is then defined the following way:
\[
{\bf ext}_u : M \mapsto M \circledast_R (R \otimes_u S)
\]

\section{The category ${\bf Mod}^{\ominus}$ of modules and retromorphisms}

Let $[M, N]$ denote the internal hom of $M$ and $N$ in $\mathcal{A}$. A \emph{module retromorphism}
\[
(u, f) : (S, N) \to (R, M)
\]
is a pair $(u, f)$ where $u : R \to S$ is a ring morphism and $f : N \to M$ is a morphism in $\mathcal{A}$ such that the following diagram commutes:
\[
\begin{tikzcd}
  R \otimes M \ar[d, "\operatorname{act}_M"'] & R \otimes N \ar[l, "R \otimes f"'] \ar[r, "u \otimes N"] & S \otimes N \ar[d, "\operatorname{act}_N"]\\ 
  M && N \ar[ll, "f"]
\end{tikzcd}
\]
The category ${\bf Mod}^{\ominus}$ has ring modules as objects and module retromorphisms. There is a forgetful functor $\pi : {\bf Mod}^{\ominus} \to {\bf Ring}$
transporting every retromorphism to the underlying ring morphism and $\pi$ is also a Grothendieck fibration. 
$\pi$ is also a bifibration. Let $u : R \to S$ be a ring morphism, then functor defined by coextension of scalar along $u$
\[
{\bf coext}_u : {\bf Mod}_R \to {\bf Mod}_S
\]
transport every $R$-module $(M, \operatorname{act}_M)$ to the $S$-module $[S,M]_u$ defined as the coreflexive equaliser of the below diagram:
\[
\begin{tikzcd}
  {[S,M]} \ar[rrrr, shift left=4, "{[u \otimes S, M]} \circ {[m_S, M]}"] \ar[rrrr, shift right=4, "{[R  \otimes S, \operatorname{act}_M]} \circ {[R \otimes \underline{\:\:}, R \otimes \underline{\:\:}]}"'] &&&& {[R \otimes S, M]} \ar[llll, "{[e_R \otimes M, M]}" description]
\end{tikzcd}
\]
Note that $[S,M]_u$ provides an internal description of the set of morphisms $f : S \to M$ making the following diagram commute:
\[
\begin{tikzcd}
  R \otimes S \ar[rr, "R \otimes f"] \ar[d, "u \otimes S"'] && R \otimes M \ar[dd, "\operatorname{act}_M"]\\
  S \otimes S \ar[d, "m_S"'] \\
  S \ar[rr, "f"'] && M
\end{tikzcd}
\]
or, equivalently, as the set of morphisms $f : {\bf res}_u S \to M$. Thus, every ring morphism $u : R \to S$ induces the following sequence of 
adjunctions:
\[
\begin{tikzcd}
  {\bf Mod}_R \ar[rrr, shift left=4, "{\bf coext}_u", ""{name=coext}] \ar[rrr, shift right=4, "{\bf ext}_u"', ""{name=ext}] &&& {\bf Mod}_S \ar[lll, "{\bf res}_u" description, ""{name=res}]
  \arrow[phantom, from=coext, to=res, "\vdash" rotate=-90]
  \arrow[phantom, from=ext, to=res, "\vdash" rotate=-90]
\end{tikzcd}
\]

\section{The category ${\bf Mod}_R$ is symmetric monoidal closed}

It is well known that the category ${\bf Mod}_R$ of modules over a ring $R$ is symmetric monoidal closed. Its tensor product
\[
\otimes_R : {\bf Mod}_R \times {\bf Mod}_R \to {\bf Mod}_R
\]
transports every pair of $R$-modules $(M,N)$ into the coreflexive equaliser of the following diagram calculated in $\mathcal{A}$:
\[
\begin{tikzcd}
  M \otimes R \otimes N \ar[rrr, shift left=4, "\operatorname{act}_M \otimes N"] \ar[rrr, shift right=4, "M \otimes \operatorname{act}_N"'] &&& M \otimes N \ar[lll, "M \otimes e_R \otimes N" description]
\end{tikzcd}
\]
In ${\bf Mod}_R$, the tensorial unit is the $R$ itself seen as an $R$-module. The internal hom
\[
  [.,.]_R : {\bf Mod}^{op}_R \times {\bf Mod}_R \to {\bf Mod}_R
\]
takes every pair of $R$-modules $M$ and $N$ to the coreflexive equaliser of the following diagram:
\[
\begin{tikzcd}
  {[M,N]} \ar[rrrr, shift left=4, "{[\operatorname{act}_M, N]}"] \ar[rrrr, shift right=4, "{[R \otimes M,\operatorname{act}_N]} \circ {[R \otimes \underline{\:\:},R \otimes \underline{\:\:}]}"'] &&&& {[R \otimes M, N]} \ar[llll, "{[e_R \otimes M, N]}" description]
\end{tikzcd}
\]

\section{The category ${\bf Mod}$ as a symmetric monoidal ringed category}

Let us consider the category ${\bf Mod} \times_{\bf Ring} {\bf Mod}$ defined by the pullback diagram as follows:
\[
\begin{tikzcd}
  {\bf Mod} \times_{\bf Ring} {\bf Mod} \ar[d] \ar[rr] && {\bf Mod} \ar[d, "\pi"] \\
  {\bf Mod} \ar[rr, "\pi"'] && {\bf Ring}
\end{tikzcd}
\]
The fibrewise tensor product 
\[
\otimes_{\bf Mod} : {\bf Mod} \times_{\bf Ring} {\bf Mod} \to {\bf Mod}
\]
transports every pair of modules $(R, M)$ and $(R, N)$ on the same monoid $R$ to the $R$-module $(R, M \otimes_R N)$, so every pair
of module homomorphisms:
\begin{center}
  $u : R \to S \models h_1 : M_1 \to N_1$

  $u : R \to S \models h_2 : M_2 \to N_2$
\end{center}
above the same ring morphism $u : R \to S$ to the module homomorphism:
\[
u : R \to S \models h_1 \otimes_u h_2 : M_1 \otimes_R M_2 \to N_1 \otimes_R N_2
\]
where $h_1 \otimes_u h_2$ is a unique morphism making the following diagram commute:
\[
\begin{tikzcd}
  M_1 \otimes R \otimes M_2 \ar[d, shift left=2, "\operatorname{act}_{M_1} \otimes M_2"] \ar[d, shift right=2, "M_1 \otimes \operatorname{act}_{M_2}"'] \ar[rrr, "h_1 \otimes u \otimes h_2"] &&& N_1 \otimes S \otimes N_2 \ar[d, shift left=2, "\operatorname{act}_{N_1} \otimes N_2"] \ar[d, shift right=2, "N_1 \operatorname{act}_{N_2} "']  \\
  M_1 \otimes M_2 \ar[d, "quotient"'] \ar[rrr, "h_1 \otimes h_2"] &&& N_1 \otimes N_2 \ar[d, "quotient"] \\
  M_1 \otimes_R M_2 \ar[rrr, "h_1 \otimes_u h_2"'] &&& N_1 \otimes_S N_2
\end{tikzcd}
\]
The fibrewise unit $\mathds{1}_{\bf Mod} : {\bf Ring} \to {\bf Mod}$ takes every commutative ring and converts it to itself viewed as a module.
Categorically, one can understood that with the concept of a \emph{ringed category}, which defined as a pair $(\mathcal{C}, \pi)$
where $\mathcal{C}$ is a category and $\pi : \mathcal{C} \to {\bf Ring}$ is a functor. The slice 2-category $\mathcal{C}/{\bf Ring}$
consists of ringed categories, fibrewise functions as 1-cells and natural transformations as 2-cells. The 2-category $\mathcal{C}/{\bf Ring}$
is Cartesian; the fibrewise product and the unit define a symmetric pseudomonoid structure on the ringed category $({\bf Mod}, \pi)$
in $\mathcal{C}/{\bf Ring}$.

\section{The categories ${\bf Alg}$ abd ${\bf CoAlg}$ of (co)commutative (co)algebras}

Let $R$ be a commutative ring in $\mathcal{A}$, a \emph{commutative $R$-algebra} is a commutative monoid in the symmetric monoidal category
${\bf Mod}_R$. A \emph{commutative algebra} is a pair $(R, A)$ where $R$ is a commutative ring and $A$ is a commutative $R$-algebra.
An algebra morphism $(u, f) : (R, A) \to (S, B)$ is a pair $(u, f)$ where $u$ is a ring morphism and $f : A \to B$ is a module morphism such that the 
following diagrams commute:
\[
\begin{tikzcd}
  A \otimes_R A \ar[d, "m_A"'] \ar[rr, "f \otimes_u f"] && B \otimes_S B \ar[d, "m_B"] & 1_R \ar[r, "u"] \ar[d, "e_A"'] & 1_S \ar[d, "e_B"] \\
  A \ar[rr, "f"'] && B                                                                 & A \ar[r, "f"'] & B
\end{tikzcd}
\]
Assume that for every commutative ring $R$ the symmetric monoidal category ${\bf Mod}_R$ has free commutative monoids, which means that the forgetful 
functor
\[
{\bf Forget}_R : {\bf Alg}_R \to {\bf Mod}_R
\]
has a left adjoint
\[
{\bf Sym}_R : {\bf Mod}_R \to {\bf Alg}_R
\]
and the adjunction ${\bf Mod} \dashv {\bf Alg}$ is also fibrewise above ${\bf Ring}$.

Given a commutative ring $R$, a \emph{cocommutative $R$-coalgebra $K$} is defined as a cocommutative comonoid in the symmetric monoidal category
${\bf Mod}_R$. Let us also assume that for any commutative ring $R$, ${\bf Mod}_R$ has cofree cocommutative comonoids, then the forgetful functor
\[
{\bf Forget}_R : {\bf CoAlg}_R \to {\bf Mod}_R
\]
have a right adjoint  
has a left adjoint
\[
{\bf Sym}_R : {\bf Mod}_R \to {\bf CoAlg}_R
\]
 
\bibliographystyle{alpha} 
\bibliography{Text}
\end{document}